\documentclass[11pt]{article}

\usepackage[spanish,es-nodecimaldot]{babel}
\usepackage[utf8]{inputenc}
\usepackage[T1]{fontenc}
\usepackage[a4paper,margin=2.5cm]{geometry}
\usepackage{graphicx}
\usepackage{booktabs}
\usepackage{array}
\usepackage{amsmath}
\usepackage{amssymb}
\usepackage{xcolor}
\usepackage{hyperref}
\usepackage{float}
\usepackage{enumitem}
\usepackage{caption}

\geometry{margin=2.5cm}
\hypersetup{colorlinks=true,linkcolor=black,citecolor=black,urlcolor=blue}
\setlist[itemize]{topsep=0.25em,itemsep=0.15em}
\setlist[enumerate]{topsep=0.25em,itemsep=0.15em}

% Rutas esperadas dentro del repositorio. Si el documento se compila desde
% latex/observational_shadow/, estas rutas apuntan a los outputs del subproyecto.
\newcommand{\figroot}{../../outputs/observational_shadow/figures}
\newcommand{\tabroot}{../../outputs/observational_shadow/tables}

% Insercion robusta de figuras: permite compilar el LaTeX aunque las figuras
% no esten presentes en la maquina actual.
\newcommand{\maybeincludegraphics}[2][]{%
  \IfFileExists{#2}{%
    \includegraphics[#1]{#2}%
  }{%
    \fbox{\parbox{0.88\linewidth}{\centering
    Figura no disponible en la ruta esperada:\\[0.4em]
    \texttt{\detokenize{#2}}\\[0.4em]
    Copiar los PDF generados por \texttt{observational\_shadow} o ajustar \texttt{\string\figroot}.}}
  }%
}

\title{Comunidades topológicas e incompletitud observacional en el catálogo de exoplanetas}
\author{Proyecto Mapper/TDA de exoplanetas}
\date{\today}

\begin{document}
\maketitle

\begin{abstract}
Este reporte interpreta los resultados del subproyecto de \emph{sombra observacional} aplicado a comunidades Mapper del catálogo de exoplanetas. La motivación no es repetir que los métodos de descubrimiento tienen sesgos, sino usar esos sesgos como señal diagnóstica para localizar regiones topológicas posiblemente submuestreadas. El resultado principal es que las sombras más claras aparecen como comunidades localizadas, con contraste fuerte frente a sus vecinos, baja o moderada imputación y dominancia de velocidad radial en los candidatos superiores. La lectura científica prudente es que estas comunidades no prueban planetas faltantes, pero sí priorizan vecindarios físico-orbitales donde el catálogo puede estar incompleto hacia planetas de menor masa, menor señal radial o menor detectabilidad instrumental.
\end{abstract}

\section{Resumen ejecutivo}

El análisis previo de auditoría observacional mostró que la topología Mapper no es neutral respecto al método de descubrimiento. Este reporte da un paso adicional: usa la composición local de cada nodo y la compara con la de sus vecinos topológicos para estimar un índice de \emph{sombra observacional}. La idea central es:
\[
\text{sesgo por método} \;\longrightarrow\; \text{frontera topológica local} \;\longrightarrow\; \text{candidato a incompletitud observacional}.
\]

El resultado más importante es que los principales candidatos de sombra observacional están dominados por \textbf{velocidad radial}. Por tanto, la interpretación más natural no es ``faltan planetas en general'', sino una hipótesis más específica:
\[
\boxed{\text{las comunidades RV de alta sombra son candidatas a incompletitud hacia masas menores o señales radiales débiles.}}
\]

Esta conclusión es propositiva pero no definitiva. El análisis no estima cuántos planetas faltan ni confirma poblaciones no observadas. Para eso se necesitarían funciones de completitud instrumental, modelos de detectabilidad o simulaciones de inyección-recuperación.

\section{Motivación astronómica}

Los sesgos por método de descubrimiento son conocidos. El método de tránsito favorece periodos cortos, radios detectables y alineación geométrica favorable. La velocidad radial favorece planetas suficientemente masivos o configuraciones que produzcan señales radiales fuertes. La imagen directa favorece planetas grandes, jóvenes, calientes y separados de su estrella. Microlensing corresponde a una ventana geométrica particular y no debe interpretarse por sí sola como una clase física.

La contribución de este subproyecto no es descubrir esos sesgos, sino preguntar si aparecen como geometría local dentro del Mapper. En otras palabras, el objetivo es convertir una limitación observacional en una herramienta de priorización:
\[
\text{método dominante} + \text{contraste con vecinos} + \text{baja imputación}
\Rightarrow
\text{región candidata a submuestreo}.
\]

\begin{table}[H]
\centering
\small
\caption{Direcciones esperadas de incompletitud por método dominante.}
\label{tab:direcciones}
\begin{tabular}{p{0.24\linewidth}p{0.66\linewidth}}
\toprule
\textbf{Método dominante} & \textbf{Dirección plausible de incompletitud} \\
\midrule
Transit & Planetas de menor radio, periodos más largos o menor probabilidad geométrica de tránsito. \\
Radial Velocity & Planetas de menor masa, menor amplitud RV, señales dinámicas más débiles o mayor ruido estelar/instrumental. \\
Imaging & Planetas menos luminosos, menos masivos, más viejos o con menor separación angular respecto a la estrella. \\
Microlensing & Regiones dependientes de geometrías de lente favorables; no debe leerse como taxonomía física directa. \\
\bottomrule
\end{tabular}
\end{table}

\section{De la auditoría de sesgo a la sombra observacional}

La auditoría previa respondía una pregunta global: si los nodos Mapper están asociados con el método de descubrimiento. La sombra observacional responde una pregunta local:
\begin{quote}
¿En qué nodos la composición por método difiere fuertemente de la composición de sus vecinos topológicos?
\end{quote}

La diferencia es importante. Un nodo puro por método no necesariamente es una sombra observacional; podría ser simplemente una región física distinta. En cambio, un nodo puro que está físicamente cerca de sus vecinos pero observacionalmente contrastado es más interesante, porque sugiere una frontera de selección.

\section{Datos y grafos analizados}

El análisis se centra en \texttt{orbital\_pca2\_cubes10\_overlap0p35}. Esta configuración es la más adecuada para la interpretación principal porque combina complejidad topológica con baja dependencia de imputación. Cuando existen artefactos compatibles, se resumen también:
\begin{itemize}
  \item \texttt{phys\_min\_pca2\_cubes10\_overlap0p35};
  \item \texttt{joint\_no\_density\_pca2\_cubes10\_overlap0p35};
  \item \texttt{joint\_pca2\_cubes10\_overlap0p35};
  \item \texttt{thermal\_pca2\_cubes10\_overlap0p35}.
\end{itemize}

La comparación entre configuraciones debe interpretarse con cautela: un índice alto en el espacio térmico no tiene la misma fuerza que un índice alto en el espacio orbital, porque el espacio térmico depende mucho más de imputación.

\section{Metodología}

Sea $G=(V,E)$ el grafo Mapper. Cada nodo $v\in V$ contiene un conjunto de planetas. Para cada método de descubrimiento $m$, se define:
\[
p_v(m)=\frac{n_v(m)}{\sum_m n_v(m)}.
\]

El vecindario topológico de $v$ es:
\[
N(v)=\{u:(u,v)\in E\}.
\]

La composición observacional del vecindario es:
\[
p_{N(v)}(m)=\frac{\sum_{u\in N(v)} n_u(m)}{\sum_m\sum_{u\in N(v)}n_u(m)}.
\]

El contraste local por método se resume con:
\[
D_v(m)=p_{N(v)}(m)-p_v(m),\qquad
R_v(m)=\frac{p_v(m)+\epsilon}{p_{N(v)}(m)+\epsilon}.
\]

La frontera composicional se mide con distancia L1:
\[
B(v)=\mathrm{method\_l1\_boundary}(v)=\sum_m |p_v(m)-p_{N(v)}(m)|.
\]

El índice principal de sombra observacional es:
\[
\mathrm{Shadow}(v)=f_{\mathrm{top}}(v)\,[1-H_{\mathrm{norm}}(v)]\,[1-I(v)]\,B(v)\,S(v),
\]
con:
\[
S(v)=\frac{\log(1+n_v)}{\log(1+n_{\max})}.
\]

Aquí $f_{\mathrm{top}}$ es la fracción del método dominante, $H_{\mathrm{norm}}$ la entropía normalizada de métodos, $I(v)$ la fracción media de imputación, $B(v)$ el contraste local de método y $S(v)$ una ponderación por tamaño nodal. La forma multiplicativa implica que un nodo recibe sombra alta solo si combina varios criterios a la vez: dominancia metodológica, baja mezcla, baja imputación, contraste con vecinos y tamaño no trivial.

\section{Resultados e interpretación}

\subsection{Distribución espacial de la sombra observacional}

\begin{figure}[H]
\centering
\maybeincludegraphics[width=0.92\linewidth]{\figroot/orbital_graph_by_shadow_score.pdf}
\caption{Grafo orbital coloreado por \texttt{shadow\_score}. Valores altos sugieren una posible frontera de selección local.}
\label{fig:shadow-score}
\end{figure}

La Figura~\ref{fig:shadow-score} muestra que los valores altos de \texttt{shadow\_score} no se distribuyen de forma homogénea sobre el grafo. Aparecen en regiones localizadas, especialmente en zonas ramificadas del Mapper orbital. Esto es relevante porque la hipótesis de sombra observacional requiere localización: si todo el grafo tuviera valores altos, el índice no distinguiría regiones prioritarias.

La lectura principal es:
\[
\text{sombra alta localizada} \Rightarrow \text{posible frontera local de selección observacional}.
\]

\subsection{Clases heurísticas de sombra}

\begin{figure}[H]
\centering
\maybeincludegraphics[width=0.92\linewidth]{\figroot/orbital_graph_by_shadow_class.pdf}
\caption{Clasificación heurística de nodos por sombra observacional.}
\label{fig:shadow-class}
\end{figure}

La Figura~\ref{fig:shadow-class} separa nodos de alta confianza, nodos de muestra pequeña, nodos mixtos o de baja sombra y nodos sin información suficiente de vecinos. El resultado deseable es que los nodos de alta sombra sean pocos y localizados, no que todo el grafo quede marcado como problemático. En ese sentido, el índice parece funcionar como un filtro de priorización.

La categoría \texttt{small\_sample\_shadow} debe tratarse con especial cautela: un nodo pequeño puede parecer puro por azar. Por eso los mejores candidatos no son simplemente los de mayor pureza, sino los que combinan sombra alta, baja imputación y tamaño suficiente.

\subsection{Sombra observacional e imputación}

\begin{figure}[H]
\centering
\maybeincludegraphics[width=0.76\linewidth]{\figroot/orbital_shadow_vs_imputation.pdf}
\caption{Relación entre sombra observacional e imputación media.}
\label{fig:shadow-imputation}
\end{figure}

La Figura~\ref{fig:shadow-imputation} muestra que varios nodos con sombra alta tienen imputación cercana a cero. Este es uno de los resultados más favorables del análisis, porque reduce la posibilidad de que la sombra observacional sea solo un artefacto del pipeline de imputación. La interpretación prudente es:
\[
\text{sombra alta + baja imputación} \Rightarrow \text{candidato confiable a frontera observacional}.
\]

Los nodos con sombra alta e imputación mayor no deben descartarse automáticamente, pero deben clasificarse como candidatos secundarios o sensibles a imputación.

\subsection{Contraste método-vecindario}

\begin{figure}[H]
\centering
\maybeincludegraphics[width=0.76\linewidth]{\figroot/orbital_method_boundary_vs_shadow.pdf}
\caption{Sombra observacional frente al contraste local de métodos.}
\label{fig:method-boundary}
\end{figure}

La Figura~\ref{fig:method-boundary} muestra una relación fuerte entre \texttt{method\_l1\_boundary} y \texttt{shadow\_score}. Esto es esperable porque la frontera L1 forma parte del índice. Aun así, la figura es útil como diagnóstico: confirma que los nodos de mayor sombra son precisamente aquellos cuya composición por método difiere de sus vecinos.

La interpretación conceptual es:
\[
\text{nodo puro} + \text{vecindario distinto} \neq \text{simple pureza};
\]
\[
\text{nodo puro} + \text{vecindario distinto} = \text{candidato a frontera de selección}.
\]

En la figura, los puntos de mayor sombra se concentran en comunidades dominadas por velocidad radial. Esto orienta la hipótesis física hacia incompletitud en masa o amplitud radial.

\subsection{Distancia físico-orbital al vecindario}

\begin{figure}[H]
\centering
\maybeincludegraphics[width=0.76\linewidth]{\figroot/orbital_physical_neighbor_distance_vs_shadow.pdf}
\caption{Sombra observacional frente a distancia físico-orbital al vecindario.}
\label{fig:physical-distance}
\end{figure}

La Figura~\ref{fig:physical-distance} es clave para separar dos explicaciones. Si un nodo tiene sombra alta y además está muy lejos físicamente de sus vecinos, la diferencia podría ser una separación física real. En cambio, si tiene sombra alta y distancia físico-orbital baja o moderada, la explicación de frontera observacional gana plausibilidad.

La región más interesante de la figura es, por tanto:
\[
\boxed{\text{sombra alta} + \text{distancia física baja/moderada}.}
\]

Esa combinación sugiere que el nodo no es completamente distinto de su entorno en variables físicas, pero sí está separado por método de descubrimiento. Es el caso más cercano a la idea de comunidad topológica submuestreada.

\subsection{Comparación entre configuraciones}

\begin{figure}[H]
\centering
\maybeincludegraphics[width=0.82\linewidth]{\figroot/config_shadow_comparison.pdf}
\caption{Comparación resumida entre configuraciones disponibles.}
\label{fig:config-comparison}
\end{figure}

La Figura~\ref{fig:config-comparison} indica que la sombra observacional aparece en varias configuraciones, no solo en el Mapper orbital. Sin embargo, la interpretación más confiable sigue siendo la orbital. El espacio térmico puede mostrar nodos de sombra alta, pero su lectura científica es más débil por alta imputación. El espacio conjunto es útil como sensibilidad, aunque mezcla variables orbitales, físicas y derivadas.

La jerarquía interpretativa recomendada es:
\[
\text{orbital} > \text{joint / joint no density} > \text{thermal}.
\]

\section{Comunidades candidatas a incompletitud}

\begin{figure}[H]
\centering
\maybeincludegraphics[width=0.86\linewidth]{\figroot/orbital_top_shadow_candidates.pdf}
\caption{Comunidades principales por índice de sombra observacional.}
\label{fig:top-candidates}
\end{figure}

La Figura~\ref{fig:top-candidates} contiene el resultado sustantivo más importante del subproyecto: las comunidades superiores por \texttt{shadow\_score} están dominadas por \textbf{Radial Velocity}. Esto desplaza la interpretación desde una afirmación genérica de sesgo hacia una hipótesis más específica:
\[
\boxed{\text{las comunidades RV de alta sombra son candidatos a incompletitud hacia planetas menos masivos.}}
\]

En términos físicos, velocidad radial detecta mejor planetas que inducen señales dinámicas suficientemente fuertes en su estrella. Por tanto, una comunidad topológica dominada por RV, contrastante con su vecindario y de baja imputación, puede interpretarse como una región donde el catálogo observa preferentemente la parte de alta masa o alta señal de una población más amplia.

La tabla \texttt{top\_shadow\_candidates.csv} debe leerse como lista de prioridades, no como lista de descubrimientos. Cada fila representa una comunidad candidata que debería inspeccionarse con variables como masa, periodo, semieje mayor, tipo de estrella, ruido instrumental y disponibilidad de seguimiento por tránsito o RV.

\section{Conclusiones por nivel de fuerza}

\subsection*{Nivel 1: conclusión segura}

Se puede afirmar que el índice de sombra identifica comunidades localizadas con alta dominancia metodológica y fuerte contraste contra sus vecinos. Estas comunidades no están distribuidas al azar sobre el grafo.

\subsection*{Nivel 2: conclusión plausible}

Es plausible interpretar las comunidades RV de alta sombra como vecindarios físico-orbitales donde el catálogo puede estar incompleto hacia objetos de menor masa o menor amplitud radial. Esta afirmación es consistente con la función de selección esperada de velocidad radial.

\subsection*{Nivel 3: conclusión no demostrada}

No se puede afirmar que esas comunidades contengan planetas faltantes confirmados, ni estimar cuántos objetos faltan, ni convertirlas en clases planetarias reales. Para eso se requieren modelos de completitud o datos instrumentales adicionales.

\section{Discusión}

La principal ganancia conceptual del subproyecto es que cambia la función del sesgo. En la auditoría observacional, el sesgo era una advertencia: impedía leer Mapper como taxonomía física directa. En la sombra observacional, el sesgo se vuelve una herramienta: permite priorizar regiones donde el catálogo puede estar incompleto.

La síntesis es:
\[
G_{\mathrm{Mapper}} = \text{estructura física} + \text{función de selección} + \text{imputación} + \text{decisiones del pipeline}.
\]

El índice de sombra intenta localizar términos donde la función de selección domina de forma local, pero sin perder trazabilidad de imputación ni tamaño nodal.

\section{Limitaciones}

\begin{enumerate}
  \item Mapper no prueba planetas faltantes; solo localiza regiones candidatas a incompletitud.
  \item Un vecino topológico no equivale a un planeta real no detectado ni a un planeta del mismo sistema estelar.
  \item Sin funciones de completitud instrumental no se pueden estimar cantidades absolutas de planetas faltantes.
  \item Los nodos pequeños pueden tener pureza artificialmente alta por tamaño muestral.
  \item La imputación puede inflar o distorsionar algunas regiones, aunque los mejores candidatos orbitales parecen tener baja imputación.
  \item El índice de sombra es heurístico y multiplicativo; debe leerse como priorización, no como probabilidad calibrada.
\end{enumerate}

\section{Trabajo futuro}

La continuación natural tiene tres niveles.

\subsection{Completitud por método}

Incorporar funciones aproximadas de detectabilidad:
\[
P(\mathrm{detectado}\mid R_p,M_p,P,a,\mathrm{estrella},m),
\]
para traducir sombra observacional en estimaciones calibradas de incompletitud.

\subsection{Validación contrafactual}

Construir Mappers balanceados por método o residualizar variables contra \texttt{discoverymethod}, \texttt{disc\_year} y \texttt{disc\_facility}. El objetivo sería medir qué parte del grafo sobrevive cuando se controla explícitamente el método de descubrimiento.

\subsection{Inspección astrofísica de comunidades RV}

Para las comunidades RV de mayor sombra, revisar masa, periodo, semieje mayor, tipo de estrella, amplitud radial esperada y posibilidad de seguimiento por tránsito o campañas RV de mayor precisión.

\section{Conclusión general}

El subproyecto confirma que Mapper puede ir más allá de la auditoría de sesgo. Las comunidades de alta sombra, especialmente las dominadas por velocidad radial, funcionan como candidatos topológicos a incompletitud observacional. El resultado no confirma planetas faltantes, pero sí identifica vecindarios donde futuras observaciones o modelos de completitud deberían concentrarse.

\[
\boxed{\text{Mapper no solo mapea planetas observados; también puede señalar dónde el catálogo deja sombras.}}
\]

\end{document}
